\section{结论}\label{sec:conclusion}
本文将 NPU 核函数内存规划重构为两个可分离、可审计的问题：哪个合法拓扑序暴露出可处理的驻留前沿，以及容量受限时应保持哪些有备份或生成的驻留间隙。生产层面的答案是围绕依赖前沿序、按真实目标选择的有界组合；固定序层面的研究答案是加权间隙下界与连续地址流量证书。恒等式 $\mathrm{Tr}=\mathrm{Vol}+\mathrm{Dt}$ 使每次 P2 变化都可分解审计。

证据有效但边界明确：公开实例上 P2 于五例改进、一例持平，P3 时间于五例改进、一例回退；两个非统一修复个案改进两个卷积实例；经标准验证的合成套件复评确立非回归与小规模同序流量最优，而非相对前代的新泛化。因此本文的结论不是普适的 clean/dirty 调度定律，而是一座边界明确的精确—启发式桥梁，以及在所研究核函数族上经验证的改进。后续工作应在超时与验证回退的保护下集成精确求解、设计更廉价且统一的代价感知序修复，并在读写语义比当前静态标签更丰富的硬件轨迹上评估本方法。
